% Generated draft: requires mathematical and source review.
\documentclass[11pt]{article}
\usepackage[margin=25mm]{geometry}
\usepackage{fontspec}
\setmainfont{Latin Modern Roman}
\usepackage{amsmath,amssymb,mathtools,unicode-math}
\setmathfont{Latin Modern Math}
\usepackage{enumitem,needspace,xurl,hyperref}
\setlength{\parindent}{0pt}
\setlength{\parskip}{5pt}
\setlength{\emergencystretch}{3em}
\newcommand{\N}{\mathbb{N}}
\newcommand{\Z}{\mathbb{Z}}
\newcommand{\Q}{\mathbb{Q}}
\newcommand{\R}{\mathbb{R}}
\newcommand{\C}{\mathbb{C}}
\newcommand{\F}{\mathbb{F}}
\newcommand{\A}{\mathbb{A}}
\newcommand{\PP}{\mathbb{P}}
\newcommand{\E}{\mathbb{E}}
\newcommand{\eps}{\varepsilon}
\newcommand{\dd}{\,\mathrm{d}}
\begin{document}
\section*{UnsolvedMath ID: 20001099}
\textbf{Source problem code:} catalogue entry
\subsection{Definitions and mathematical statement}
Let $A$ be a finite nonempty subset of an abelian group $G$, let $D=A-A$, and count ordered representations by
\[
 r_A(d)=\#\{(a',a'')\in A^2:a'-a''=d\}.
\]
V. Lev's source records that $r_A(d)>|A|/2$ for every $d\in D$ forces $D$ to be a subgroup. At the non-strict threshold, the example $A=H\cup(g+H)$, where $H$ is a finite subgroup and $g+H$ has order at least five in $G/H$, has
$D=(-g+H)\cup H\cup(g+H)$, not a subgroup, while every represented difference has at least $|H|=|A|/2$ representations.

The source asks two separate questions: whether $r_A(d)\geq|A|/2$ for every $d\in D$ always makes $D$ a subgroup or a union of three cosets of a common subgroup, and whether the same conclusion holds under the weaker strict hypothesis
\[
 r_A(d)>|A|/3\qquad(d\in D).
\]
The strictness and the distinction between these two thresholds are essential.

The attached aid refutes the one-third assertion with $G=\Z/3\Z\times\Z/7\Z$ and
\[
 A=(\Z/3\Z\times\{0\})\cup(\{0,1\}\times\{1\})\cup(\Z/3\Z\times\{2\}).
\]
It records $|A|=8$ and $D=\Z/3\Z\times\{0,\pm1,\pm2\}$. On heights $\pm2$ the multiplicity is $3$ for each first coordinate; on heights $\pm1$ it is $4$; on height zero the values at first coordinates $0,1,2$ are respectively $8,7,7$. Thus every $d\in D$ has $r_A(d)\geq3>8/3$. The $15$-element set $D$ is neither a subgroup nor a union of at most three cosets of any common subgroup of the order-$21$ group. This is an odd-order example, so absence of nonzero two-torsion does not rescue the one-third conclusion. Its minimum relative multiplicity $3/8$ does not satisfy the half-threshold hypothesis and therefore does not settle that separate question.

More generally, in the same three-layer construction with full outer fibers of size $h$ and a nonempty middle fiber of size $m\leq h$, the aid records relative minimum $\min(h,2m)/(2h+m)$. This is at most $2/5$, with equality at $m=h/2$ when that integer size is available. This construction-specific ratio is not claimed to be a universal classification threshold.
\subsection{Short English statement}
Ask whether uniformly popular differences force a subgroup or three-coset structure. The supplied eight-element odd-order example refutes the strict one-third version. It does not answer the separate question at multiplicity at least one half.
\subsection{Sources}
\begin{itemize}
\item AIM list \emph{Recent trends in additive combinatorics}, problem 1.18 attributed to V. Lev: \url{https://aimath.org/WWN/additivecomb/additivecomb.pdf}.
\item Local registry \texttt{data/unsolvedmath/problems.json}, permanent ID 20001099, exact source code \texttt{catalogue entry}; canonical location \texttt{aim-combinatorics-notes.json notes[223]}.
\item Attached attempt 1, counterexample, problem status at run refuted, rigor proved, low novelty confidence. The explicit example and general three-layer ratio have no separate complete external reference in the supplied record; the aid is machine-generated, not peer reviewed, requiring independent verification before citation.
\end{itemize}
\end{document}
